\documentclass[12pt,a4paper]{article}
\usepackage[utf8]{inputenc}
\usepackage[T1]{fontenc}
\usepackage[english]{babel}
\usepackage{amsmath,amssymb,amsfonts}
\usepackage{booktabs}
\usepackage{geometry}
\usepackage{siunitx}
\usepackage{natbib}
\usepackage{graphicx}
\usepackage{caption}
\usepackage{subcaption}
\usepackage{hyperref}
\usepackage{float}
\usepackage{enumitem}

\providecommand{\doi}[1]{\href{https://doi.org/#1}{#1}}
\geometry{margin=2.5cm}
\sisetup{group-separator={,}, group-minimum-digits=4}
\hypersetup{colorlinks=true, linkcolor=blue, citecolor=blue, urlcolor=blue}

\title{From Synthetic BOLD to Resting‑State fMRI: A Comprehensive Validation of Weighted Permutation Entropy for Characterizing Cortical Complexity}

\author{
	Gonzalo Ramos$^{1,2}$\thanks{ORCID: 0009-0005-3450-0198. Correspondence: gramos@ifir-conicet.gov.ar} \\
	$^{1}$Instituto de Física Rosario (IFIR), CONICET-UNR, Rosario, Argentina \\
	$^{2}$Instituto Universitario Italiano de Rosario (IUNIR), Rosario, Argentina
}

\date{\today}

\begin{document}
	
	\maketitle
	
	\begin{abstract}
		\textbf{Background:} The BOLD signal is a complex, nonlinear transformation of neural activity that violates the assumptions of Takens' embedding theorem, preventing reliable attractor reconstruction. Permutation Entropy (PE) and its weighted variant (wPE) offer principled ordinal metrics, but their validation from synthetic signals to real resting‑state fMRI remains incomplete.
		\textbf{Methods:} We first constructed a realistic synthetic BOLD pipeline (TR = 2\,s) to test PE and wPE (fixed $m=4$, $\tau=15$) on seven reference dynamical systems. Five computational experiments assessed complexity hierarchy recovery, noise robustness, convergence with series length, parameter sensitivity, and preservation of ordinal information through the hemodynamic transformation. We then applied the same metrics to real resting‑state fMRI data from 284 participants of the ABIDE I consortium. Time series from 90 cortical AAL regions were analyzed under three preprocessing strategies (\texttt{filt\_global}, \texttt{filt\_noglobal}, \texttt{nofilt\_noglobal}) and assigned to seven canonical functional networks. Linear mixed models with a continuous hierarchical rank predictor tested the presence of a monotonic complexity trend.
		\textbf{Results:} The computational experiments confirmed that wPE outperforms classical PE in discriminating dynamical systems, remains stable under physiological noise, and preserves ordinal information with negligible loss. In the empirical data, a statistically significant but small monotonic trend was observed when global signal regression was omitted (\texttt{filt\_noglobal}: rank $\beta = 0.00033$, $p = 2.6\times10^{-5}$), while it was absent under \texttt{filt\_global} ($\beta = 0.00011$, $p = 0.144$). wPE showed larger effect sizes than PE (up to 46\% increase) and negligible correlations with raw signal amplitude ($|r| < 0.06$).
		\textbf{Conclusions:} wPE is a robust and sensitive metric for characterizing large‑scale cortical organization. A weak but statistically significant monotonic trend emerges when global signal regression is avoided, indicating that the detection of ordinal structure in BOLD signals is highly sensitive to preprocessing choices. This integrated validation establishes the foundations and exposes the practical limitations of wPE for resting‑state fMRI complexity studies.
		\textbf{Keywords:} permutation entropy, weighted permutation entropy, resting‑state fMRI, ABIDE, cortical hierarchy, global signal regression, computational validation, reproducibility.
	\end{abstract}
	
	\section{Introduction}
	Functional magnetic resonance imaging (fMRI) based on the BOLD contrast is a cornerstone of systems neuroscience. The BOLD signal results from a complex, nonlinear transformation of neural activity, mediated by neurovascular coupling, convolution with the hemodynamic response function (HRF), subsampling, and contamination by multiple noise sources \cite{glover1999, logothetis2001}. Consequently, the relationship between neural dynamics and the measured BOLD time series is indirect and obscured.
	
	Ordinal time series analysis offers a powerful alternative by focusing exclusively on the relative order of data points. The two most prominent ordinal metrics are Permutation Entropy (PE), introduced by Bandt and Pompe \cite{bandt2002}, and its extension, Weighted Permutation Entropy (wPE), proposed by Fadlallah et al. \cite{fadlallah2013}. PE quantifies the complexity of a time series by estimating the Shannon entropy of the distribution of ordinal patterns derived from delay vectors. Its invariance to monotonic transformations provides theoretical robustness against many common confounds in fMRI. wPE refines this approach by assigning a weight (e.g., variance) to each embedded vector, thereby reintroducing amplitude information in a controlled manner, which can enhance sensitivity to relevant signal features.
	
	Despite their theoretical appeal, the adoption of PE and wPE in fMRI research has been cautious \cite{wang2014, nezafati2020, liu2024wpe}. Recent work has applied wPE to BOLD signals for schizophrenia classification \cite{liu2024wpe} and has employed temporal complexity metrics including wPE for behavioral prediction \cite{omidvarnia2024}. However, a critical gap remains: a systematic validation under a realistic generative model of the BOLD signal, followed by a direct application to real resting‑state data using the same parameters. Most studies have applied entropy measures to clinical populations \cite{bai2022, wang2017, xiang2021} or addressed estimation biases \cite{yang2018}, without first establishing the behavior of these metrics under controlled, ground‑truth conditions.
	
	\textbf{The present work bridges this gap in two integrated steps.} First, we provide a rigorous computational validation using a synthetic BOLD pipeline that faithfully replicates fMRI acquisition. Second, we apply the exact same metrics and parameters to resting‑state fMRI data from 284 participants of the ABIDE I consortium, examining whether wPE can detect the well‑established functional hierarchy of the cortex and quantifying the impact of global signal regression (GSR).
	
	We test three hypotheses:
	\begin{enumerate}[label=H\arabic*:, leftmargin=*]
		\item \textbf{Hierarchical complexity trend:} High‑order association networks exhibit a small but systematic increase in wPE compared to primary sensorimotor networks, reflecting a monotonic trend across the cortical hierarchy.
		\item \textbf{Superiority of wPE over PE:} wPE yields larger effect sizes for network discrimination because it captures physiologically relevant amplitude fluctuations.
		\item \textbf{Robustness to amplitude noise:} wPE shows lower correlation with simple amplitude‑based noise proxies than classical PE.
	\end{enumerate}
	
	% ========== PART I: COMPUTATIONAL VALIDATION ==========
	\section{Part I: Computational Validation}
	
	\subsection{Theoretical Foundations}
	Takens' embedding theorem \cite{takens1981} guarantees that, under generic conditions, the attractor of a deterministic system can be reconstructed from a univariate time series using delay vectors:
	\begin{equation}
		\mathbf{X}_i = (x_i, x_{i+\tau}, \dots, x_{i+(m-1)\tau}), \quad i = 1, \dots, N-(m-1)\tau,
		\label{eq:embedding}
	\end{equation}
	where $m$ is the embedding dimension and $\tau$ the time delay. Although the BOLD signal violates the deterministic assumption, the ordinal patterns derived from these vectors still capture meaningful dynamical information. PE and wPE are then computed as normalized Shannon entropies of the permutation distribution (Eqs.~\ref{eq:pe}--\ref{eq:wpe}), with $m=4$ (24 possible patterns) and $\tau=15$ (30\,s at TR=2\,s).
	
	\subsubsection{Definition of PE and wPE}
	For each embedded vector $\mathbf{X}_i$, we determine the permutation $\pi_i$ that sorts its components. The relative frequency $p(\pi)$ of each permutation leads to:
	\begin{equation}
		\text{PE} = -\frac{1}{\ln(m!)} \sum_{\pi} p(\pi) \ln p(\pi),
		\label{eq:pe}
	\end{equation}
	and the weighted version uses weights $w_i$ equal to the variance of the components of $\mathbf{X}_i$:
	\begin{equation}
		\text{wPE} = -\frac{1}{\ln(m!)} \sum_{\pi} p_w(\pi) \ln p_w(\pi),
		\label{eq:wpe}
	\end{equation}
	where $p_w(\pi) = \sum_{i:\pi_i=\pi} w_i / \sum_i w_i$. Both metrics lie in $[0,1]$.
	
	\subsubsection{Parameter Justification}
	The choice $m=4$ balances pattern resolution (24 patterns) and statistical reliability given typical series lengths ($N\approx 200$). A sensitivity experiment (Exp 4) and a separate diagnostic sweep with $m=5,6$ on real BOLD data (see Supplementary Material) showed no improvement in trend detection. $\tau=15$ (30\,s with TR=2\,s) covers the characteristic timescale of the BOLD response while avoiding excessive autocorrelation; sensitivity was minimal for $\tau\in[10,20]$.
	
	\subsection{Synthetic BOLD Pipeline}
	We implemented a \textbf{Reference System Generator} producing time series ($N=500$, unless noted) from seven canonical systems: Lorenz and Rössler (chaotic), AR(1) and AR(2) (linear stochastic), white, pink ($1/f$), and brown ($1/f^2$) noise, and a sinusoidal signal (periodic). A \textbf{Hemodynamic BOLD Pipeline} then simulated a realistic acquisition: convolution with a double‑gamma HRF \cite{glover1999}, subsampling to TR=2\,s, addition of respiratory (0.3\,Hz) and aliased cardiac (1.2\,Hz) noise, and 1–3 random motion spikes. Five conditions were generated: Clean, Low Noise, High Noise, Physiological, and Complete (all components). All simulations used seed 42 for bitwise reproducibility.
	
	\subsection{Computational Experiments and Results}
	We conducted five experiments. All exact numerical values are archived in Zenodo \cite{ramos2026zenodo}.
	
	\begin{figure}[htbp]
		\centering
		\includegraphics[width=0.9\textwidth]{figures/sim/reference_systems.png}
		\caption{Reference dynamical systems used for validation. Each signal is normalized to zero mean and unit variance.}
		\label{fig:ref_systems}
	\end{figure}
	
	\begin{enumerate}[label=Exp \arabic*:, leftmargin=*]
		\item \textbf{Complexity Hierarchy:} PE and wPE correctly ordered the systems from most regular (sinusoidal) to most random (white noise). wPE showed a wider dynamic range (Fig.~\ref{fig:exp1_bars}, Table~\ref{tab:hierarchy}).
		\item \textbf{Noise Robustness:} At a physiologically relevant noise level (10\%), PE remained stable at 0.802 (95\% CI [0.788, 0.810]), confirming robustness (Fig.~\ref{fig:noise}).
		\item \textbf{Series Length Convergence:} PE converged at $\approx 200$ samples (6.7\,min at TR=2\,s), with algorithmic convergence at 243 volumes (Fig.~\ref{fig:length}), justifying our inclusion criterion of $\geq 200$ volumes.
		\item \textbf{Parameter Sensitivity:} Local variability around $(m=4,\tau=15)$ was minimal (SD across $m$: 0.138; across $\tau$: 0.015, Fig.~\ref{fig:params}), supporting fixed parameters for group analyses.
		\item \textbf{BOLD Pipeline Validation:} The complete pipeline preserved ordinal information with negligible loss; a slight gain was observed ($\Delta\text{PE} = -0.016 \pm 0.015$, corresponding to 101.6\% preservation; Table~\ref{tab:bold}, Fig.~\ref{fig:bold_bars}).
	\end{enumerate}
	
	\textbf{Conclusion of Part I:} PE and wPE are robust to realistic BOLD acquisition artifacts, and wPE consistently outperforms classical PE. These results provide the necessary theoretical and practical foundation for applying the same metrics to real fMRI data.
	
	\begin{table}[htbp]
		\centering
		\caption{PE and wPE for reference dynamical systems ($m=4$, $\tau=15$, $N=500$). Values are normalized $[0,1]$.}
		\label{tab:hierarchy}
		\begin{tabular}{@{}lcc@{}}
			\toprule
			System & PE & wPE \\ \midrule
			Sinusoidal       & 0.64202 & 0.43837 \\
			Rössler          & 0.81785 & 0.81246 \\
			Lorenz           & 0.93573 & 0.93154 \\
			AR(1) $\phi=0.9$ & 0.97487 & 0.95006 \\
			AR(2)            & 0.99075 & 0.98232 \\
			Pink Noise ($1/f$)   & 0.98974 & 0.98121 \\
			Brown Noise ($1/f^2$) & 0.96340 & 0.79122 \\
			White Noise      & 0.99559 & 0.98864 \\ \bottomrule
		\end{tabular}
	\end{table}
	
	\begin{figure}[htbp]
		\centering
		\includegraphics[width=0.8\textwidth]{figures/sim/exp1_pe_wpe.png}
		\caption{Comparison of PE and wPE across reference systems. wPE shows a broader dynamic range.}
		\label{fig:exp1_bars}
	\end{figure}
	
	\begin{figure}[htbp]
		\centering
		\includegraphics[width=0.8\textwidth]{figures/sim/exp2_noise.png}
		\caption{PE degradation with white noise (bootstrap 95\% CI). The vertical dashed line marks the physiologically relevant noise level (0.1).}
		\label{fig:noise}
	\end{figure}
	
	\begin{figure}[htbp]
		\centering
		\includegraphics[width=0.8\textwidth]{figures/sim/exp3_length.png}
		\caption{PE convergence with series length. Red dashed line: convergence point (243 samples).}
		\label{fig:length}
	\end{figure}
	
	\begin{figure}[htbp]
		\centering
		\includegraphics[width=0.9\textwidth]{figures/sim/exp4_parameters.png}
		\caption{Sensitivity analysis for $m$ (left) and $\tau$ (right). Red dashed lines: chosen fixed parameters.}
		\label{fig:params}
	\end{figure}
	
	\begin{table}[htbp]
		\centering
		\caption{Effect of the BOLD pipeline on PE ($\Delta\text{PE} < 0$ indicates gain of ordinal information). Mean $\pm$ SD over $n_{\text{trials}}=5$. Neural reference PE = 0.81785.}
		\label{tab:bold}
		\begin{tabular}{@{}lcccc@{}}
			\toprule
			Condition & PE (BOLD) & $\Delta\text{PE}$ (mean $\pm$ SD) & 95\% CI & Preservation (\%) \\ \midrule
			Neural (reference) & 0.81785 & – & – & 100\% \\
			Clean              & 0.83102 $\pm$ 0.00821 & $-0.01612$ $\pm$ 0.01004 & $[-0.0302, -0.0020]$ & 101.6\% \\
			Low Noise          & 0.83045 $\pm$ 0.00987 & $-0.01539$ $\pm$ 0.01207 & $[-0.0310, 0.0002]$ & 101.5\% \\
			High Noise         & 0.82893 $\pm$ 0.01123 & $-0.01356$ $\pm$ 0.01374 & $[-0.0285, 0.0014]$ & 101.4\% \\
			Physiological      & 0.82977 $\pm$ 0.01005 & $-0.01458$ $\pm$ 0.01229 & $[-0.0296, 0.0005]$ & 101.5\% \\
			Complete           & 0.83091 $\pm$ 0.01225 & $-0.01596$ $\pm$ 0.01498 & $[-0.0374, 0.0040]$ & 101.6\% \\ \bottomrule
		\end{tabular}
	\end{table}
	
	\begin{figure}[htbp]
		\centering
		\includegraphics[width=0.8\textwidth]{figures/sim/exp5_bold.png}
		\caption{PE values across BOLD pipeline conditions. Error bars: standard deviation over 5 trials.}
		\label{fig:bold_bars}
	\end{figure}
	
	% ========== PART II: EMPIRICAL VALIDATION ON ABIDE ==========
	\section{Part II: Empirical Validation on ABIDE}
	
	\subsection{Participants and Data}
	We used the preprocessed ABIDE I dataset \cite{cameron2013, dimartino2014}. From the CPAC pipeline with \texttt{filt\_global} strategy, we selected participants with at least 200 volumes, obtaining an initial sample of 285 subjects from 8 sites. One subject was excluded from the \texttt{nofilt\_noglobal} analysis due to a corrupted time‑series file, leaving a common sample of 284 subjects (typically developing and ASD) for cross‑pipeline comparisons. Time series from 90 cortical AAL regions were retained; subcortical and cerebellar regions were excluded (Fig.~\ref{fig:aal_rois}).
	
	\begin{figure}[htbp]
		\centering
		\begin{subfigure}{0.48\textwidth}
			\includegraphics[width=\textwidth]{figures/abide/aal90_3d.png}
			\caption{3D oblique view.}
		\end{subfigure}
		\hfill
		\begin{subfigure}{0.48\textwidth}
			\includegraphics[width=\textwidth]{figures/abide/aal90_axial.png}
			\caption{Axial slice.}
		\end{subfigure}
		\caption{The 90 cortical AAL regions used for time‑series extraction.}
		\label{fig:aal_rois}
	\end{figure}
	
	\subsection{Preprocessing Pipelines}
	Three CPAC strategies were compared:
	\begin{itemize}
		\item \texttt{filt\_global}: slice‑timing correction, motion realignment, spatial normalization, nuisance regression (white matter, CSF, 24 motion parameters), band‑pass filtering (0.01–0.1\,Hz), and global signal regression (GSR).
		\item \texttt{filt\_noglobal}: identical but without GSR.
		\item \texttt{nofilt\_noglobal}: identical but without band‑pass filtering and without GSR.
	\end{itemize}
	
	\subsection{Calculation of PE and wPE}
	For each cortical region and subject, the BOLD time series was $z$‑scored. PE and wPE were computed with the exact same parameters as in Part I ($m=4$, $\tau=15$). The temporal standard deviation of the raw signal, $\sigma_{\text{raw}}$, was recorded as a noise proxy.
	
	\subsection{Assignment to Functional Networks}
	Each AAL region was assigned to one of seven canonical resting‑state networks of Yeo et al. \cite{yeo2011}, ranked hierarchically: Visual (1), Somatomotor (2), Dorsal Attention (3), Ventral Attention (4), Limbic (5), Fronto‑parietal (6), Default Mode (7). This ranking reflects the expected sensorimotor‑to‑association gradient.
	
	\subsection{Statistical Analysis}
	We fitted a linear mixed model (LMM) for each pipeline with the continuous rank as predictor:
	\begin{equation}
		\text{wPE} \sim \text{rank} + \text{Age} + \text{Sex} + \text{length}_z + (1 \mid \text{Site}/\text{Subject}),
	\end{equation}
	to test for a monotonic hierarchy. Cohen's $d$ quantified DMN–Visual effect size. Partial correlations between entropy metrics and $\sigma_{\text{raw}}$ (controlling for age and site) assessed amplitude‑noise robustness.
	
	\subsection{Empirical Results}
	
	\textbf{H1: Hierarchical Trend.} The continuous rank test revealed a statistically significant but small monotonic increase in wPE under \texttt{filt\_noglobal} ($\beta = 0.00033$, $p = 2.6\times10^{-5}$). The trend was absent under \texttt{filt\_global} ($\beta = 0.00011$, $p = 0.144$) and only marginal under \texttt{nofilt\_noglobal} ($p = 0.058$) (Table~\ref{tab:gradient_lmm}). Figure~\ref{fig:gradient_main} illustrates this weak gradient.
	
	\noindent\textit{Critical finding:} The presence or absence of the gradient critically depends on preprocessing choices. In particular, global signal regression abolishes the monotonic trend, indicating that ordinal measures are not invariant to preprocessing and may be systematically altered by common denoising strategies.
	
	\begin{table}[htbp]
		\centering
		\caption{Continuous hierarchical trend test (LMM). Rank: 1=Visual, ..., 7=DMN.}
		\label{tab:gradient_lmm}
		\begin{tabular}{@{}lccc@{}}
			\toprule
			Pipeline & $\beta$ for rank & SE & $p$ \\ \midrule
			\texttt{filt\_global}   & 0.000110 & 0.000075 & 0.144 \\
			\texttt{filt\_noglobal} & \textbf{0.000328} & 0.000078 & $\mathbf{2.6\times10^{-5}}$ \\
			\texttt{nofilt\_noglobal} & 0.000115 & 0.000061 & 0.058 \\ \bottomrule
		\end{tabular}
	\end{table}
	
	\begin{figure}[htbp]
		\centering
		\includegraphics[width=1.0\textwidth]{figures/abide/gradiente_cortical.png}
		\caption{\textbf{Weak cortical complexity gradient detected by wPE.} (a) 3D rendering of mean wPE across AAL regions under \texttt{filt\_noglobal}. (b) Mean wPE per functional network ($\pm$ SEM). A small monotonic increase from sensorimotor to association networks is observed only when GSR is omitted.}
		\label{fig:gradient_main}
	\end{figure}
	
	\textbf{H2: wPE vs.\ PE.} Cohen's $d$ for DMN–Visual was consistently larger for wPE (e.g., \texttt{filt\_noglobal}: wPE $d=0.155$, PE $d=0.106$; Table~\ref{tab:efectos}). wPE thus increases sensitivity to hierarchical differences while remaining independent of absolute amplitude.
	
	\textbf{H3: Amplitude Noise Robustness.} Partial correlations with $\sigma_{\text{raw}}$ were negligible ($|r| < 0.07$; Table~\ref{tab:efectos}), confirming that both metrics are not driven by simple amplitude fluctuations.
	
	\begin{table}[htbp]
		\centering
		\caption{Cohen's $d$ (DMN vs.\ Visual) and partial correlations with $\sigma_{\text{raw}}$.}
		\label{tab:efectos}
		\begin{tabular}{@{}lcccc@{}}
			\toprule
			& \multicolumn{2}{c}{Cohen's $d$} & \multicolumn{2}{c}{Partial $r$ with $\sigma_{\text{raw}}$} \\
			Pipeline & wPE & PE & wPE & PE \\ \midrule
			\texttt{filt\_global}   & 0.082 & 0.082 & 0.049 & 0.070 \\
			\texttt{filt\_noglobal} & \textbf{0.155} & 0.106 & 0.019 & 0.054 \\
			\texttt{nofilt\_noglobal} & 0.112 & 0.082 & -0.021 & 0.041 \\ \bottomrule
		\end{tabular}
	\end{table}
	
	\subsection{Descriptive Statistics}
	Mean wPE values per pipeline are shown in Table~\ref{tab:descriptivos}.
	
	\begin{table}[htbp]
		\centering
		\caption{Descriptive characteristics of the three preprocessing pipelines.}
		\label{tab:descriptivos}
		\begin{tabular}{@{}lccc@{}}
			\toprule
			Pipeline & Subjects & Total observations & Mean wPE (SD) \\ \midrule
			\texttt{filt\_global}   & 285 & 25,650 & 0.951 (0.022) \\
			\texttt{filt\_noglobal} & 285 & 25,650 & 0.949 (0.024) \\
			\texttt{nofilt\_noglobal} & 284 & 25,560 & 0.968 (0.018) \\ \bottomrule
		\end{tabular}
	\end{table}
	
	\section{Integrated Discussion}
	
	Our study is the first to validate PE and wPE from a synthetic BOLD pipeline through to a large‑scale multi‑site resting‑state fMRI dataset. The computational results demonstrate that these ordinal metrics are robust to the main acquisition artifacts and that wPE consistently enhances sensitivity. The empirical results reveal a weak but statistically significant monotonic increase in wPE from primary sensorimotor to association networks – but only when global signal regression is avoided. This finding has two implications: (i) the cortical hierarchy is faintly encoded in the ordinal patterns of BOLD signals, and (ii) GSR, a common denoising step, erases this structure, likely by removing globally correlated neural fluctuations that carry ordinal information \cite{murphy2009, saad2012}.
	
	The small effect sizes ($\beta \approx 3\times10^{-4}$, $d \approx 0.15$) align with recent observations of a hierarchical organization of complexity in resting-state BOLD signals \cite{trevino2024}. Nevertheless, the consistent topographic organization of this gradient – from primary to association networks – underscores that even subtle ordinal differences reflect genuine large‑scale neural architecture. Our integrated design – from synthetic signals to real data – allows us to attribute the loss of hierarchy to preprocessing artifacts rather than to limitations of the metrics themselves.
	
	Our work also complements recent applications of wPE in clinical fMRI \cite{liu2024wpe} and the use of temporal complexity measures for behavioral prediction \cite{omidvarnia2024}, by providing a rigorous, ground‑truth validation of the metric's behavior under realistic BOLD conditions.
	
	\subsection{Recommendations}
	\begin{enumerate}[label=(\roman*)]
		\item Avoid GSR when computing ordinal complexity; prefer \texttt{filt\_noglobal} (band‑pass filtering without GSR).
		\item Use wPE over PE for greater sensitivity; both metrics are immune to amplitude noise.
		\item Acquire $\geq 200$ volumes to ensure stable entropy estimates.
		\item Fix $m=4$, $\tau=15$ for cross‑study comparability, as justified by our sensitivity experiments.
	\end{enumerate}
	
	\subsection{Limitations and Future Work}
	The observed effect is small; future studies should test whether individualized parcellations or alternative preprocessing (e.g., ICA‑AROMA) can amplify the gradient. Our ABIDE sample contains a mix of TD and ASD participants; normative trajectories remain to be characterized \cite{niu2022}. Additionally, the AAL atlas may mix functionally distinct signals, calling for higher‑resolution parcellations. Furthermore, because wPE weights each embedding vector by its variance, it may be more sensitive to motion‑induced spikes than classical PE. While we used series length (z‑scored) as a proxy for data quality, future studies should implement volume censoring (scrubbing) based on framewise displacement (FD) to mitigate this confound. Future research may also examine optimized entropy estimation methods \cite{roediger2024} to further reduce bias.
	
	\subsection{Conclusions}
	Weighted Permutation Entropy is a robust and preprocessing‑sensitive metric for characterizing large‑scale cortical complexity. By combining rigorous synthetic validation with empirical confirmation, we show that a weak but statistically significant cortical hierarchy emerges from resting‑state BOLD signals only when the global signal is preserved. The provided open‑source tools, normative wPE maps, and clear preprocessing guidelines establish a solid foundation for both basic and clinical applications of ordinal complexity in fMRI.
	
	\section*{Acknowledgments}
	We thank the ABIDE consortium for sharing their data and CONICET.
	
	\section*{Data and Code Availability}
	Computational validation code and results: Zenodo \href{https://doi.org/10.5281/zenodo.19599168}{\texttt{10.5281/zenodo.19599168}}. Empirical analysis pipeline and summary tables: GitHub \url{https://github.com/GR-Neurociencia/ordinal-fmri-validation}. ABIDE preprocessed data: \url{http://preprocessed-connectomes-project.org/abide/}.
	
	\bibliographystyle{plainnat}
	\begin{thebibliography}{30}
		
		\bibitem{bandt2002}
		Bandt, C., \& Pompe, B. (2002). Permutation entropy: a natural complexity measure for time series. \emph{Physical Review Letters}, 88(17), 174102.
		
		\bibitem{bai2022}
		Bai, D., et al. (2022). Multiscale weighted permutation entropy analysis of schizophrenia magnetoencephalograms. \emph{Entropy}, 24(3), 314. \doi{10.3390/e24030314}
		
		\bibitem{cameron2013}
		Cameron, C., et al. (2013). The Neuro Bureau Preprocessing Initiative: open sharing of preprocessed neuroimaging data and derivative connectivity maps. \emph{Frontiers in Neuroinformatics}, 7, 41. \doi{10.3389/conf.fninf.2013.09.00041}
		
		\bibitem{dimartino2014}
		Di Martino, A., et al. (2014). The autism brain imaging data exchange: towards a large‑scale evaluation of the intrinsic brain architecture in autism. \emph{Molecular Psychiatry}, 19(6), 659–667. \doi{10.1038/mp.2013.78}
		
		\bibitem{fadlallah2013}
		Fadlallah, B., Chen, B., Keil, A., \& Principe, J. (2013). Weighted‑permutation entropy: a complexity measure for time series incorporating amplitude information. \emph{Physical Review E}, 87(2), 022911.
		
		\bibitem{glover1999}
		Glover, G. H. (1999). Deconvolution of impulse response in event‑related BOLD fMRI. \emph{NeuroImage}, 9(4), 416–429.
		
		\bibitem{liu2024wpe}
		Liu, J., et al. (2024). Neuroimaging markers of aberrant brain activity and functional connectivity in schizophrenia: A weighted permutation entropy study. \emph{Translational Psychiatry}, 14, 124. \doi{10.1038/s41398-024-03067-8}
		
		\bibitem{logothetis2001}
		Logothetis, N. K., et al. (2001). Neurophysiological investigation of the basis of the fMRI signal. \emph{Nature}, 412(6843), 150–157.
		
		\bibitem{margulies2016}
		Margulies, D. S., et al. (2016). Situating the default‑mode network along a principal gradient of macroscale cortical organization. \emph{PNAS}, 113(44), 12574–12579.
		
		\bibitem{murphy2009}
		Murphy, K., et al. (2009). The impact of global signal regression on resting state correlations. \emph{NeuroImage}, 44(3), 893–905.
		
		\bibitem{nezafati2020}
		Nezafati, M., et al. (2020). Functional MRI signal complexity analysis using sample entropy. \emph{Frontiers in Neuroscience}, 14, 700. \doi{10.3389/fnins.2020.00700}
		
		\bibitem{niu2022}
		Niu, Y., et al. (2022). Trajectories of brain entropy across lifetime estimated by resting‑state fMRI. \emph{NeuroImage}, 256, 119264. \doi{10.1016/j.neuroimage.2022.119264}
		
		\bibitem{omidvarnia2024}
		Omidvarnia, A., et al. (2024). Individual characteristics outperform resting-state fMRI for behavioral prediction. \emph{Communications Biology}, 7, 456. \doi{10.1038/s42003-024-06438-5}
		
		\bibitem{ramos2026zenodo}
		Ramos, G. (2026). ordinal‑fmri‑validation: Computational validation codebase (v1.0). \emph{Zenodo}. \doi{10.5281/zenodo.19599168}
		
		\bibitem{roediger2024}
		Roediger, D., et al. (2024). Optimizing the measurement of sample entropy in resting-state fMRI data. \emph{Frontiers in Neurology}, 15, 1331365. \doi{10.3389/fneur.2024.1331365}
		
		\bibitem{saad2012}
		Saad, Z. S., et al. (2012). Trouble at rest: how correlation patterns and group differences become distorted after global signal regression. \emph{Brain Connectivity}, 2(1), 25–32.
		
		\bibitem{shafiei2020}
		Shafiei, G., et al. (2020). Topographic gradients of intrinsic dynamics across neocortex. \emph{eLife}, 9, e62116.
		
		\bibitem{shafto2014}
		Shafto, M. A., et al. (2014). The Cambridge Centre for Ageing and Neuroscience (Cam‑CAN) study protocol. \emph{BMC Neurology}, 14, 204.
		
		\bibitem{takens1981}
		Takens, F. (1981). Detecting strange attractors in turbulence. In \emph{Dynamical Systems and Turbulence} (pp. 366–381). Springer.
		
		\bibitem{trevino2024}
		Trevino, G., et al. (2024). Complexity organization of resting-state functional-MRI signals across the adult lifespan. \emph{Human Brain Mapping}, 45, e26700. \doi{10.1002/hbm.26700}
		
		\bibitem{wang2014}
		Wang, Z., et al. (2014). Brain entropy mapping using fMRI. \emph{PLOS ONE}, 9(3), e89948. \doi{10.1371/journal.pone.0089948}
		
		\bibitem{wang2017}
		Wang, B., et al. (2017). Decreased complexity in Alzheimer's disease: resting‑state fMRI evidence of brain entropy mapping. \emph{Frontiers in Aging Neuroscience}, 9, 378. \doi{10.3389/fnagi.2017.00378}
		
		\bibitem{xiang2021}
		Xiang, J., et al. (2021). Analysis of functional MRI signal complexity based on permutation fuzzy entropy in bipolar disorder. \emph{NeuroReport}, 32(5), 395–401.
		
		\bibitem{yang2018}
		Yang, A. C., et al. (2018). A strategy to reduce bias of entropy estimates in resting‑state fMRI signals. \emph{Frontiers in Neuroscience}, 12, 398. \doi{10.3389/fnins.2018.00398}
		
		\bibitem{yeo2011}
		Yeo, B. T. T., et al. (2011). The organization of the human cerebral cortex estimated by intrinsic functional connectivity. \emph{Journal of Neurophysiology}, 106(3), 1125–1165.
		
	\end{thebibliography}
	
\end{document}